% ============================================================
\section{4.1 kNN: 알고리즘, 거리함수, 편향-분산 트레이드오프}
\begin{frame}{``가장 가까운 5명의 평균''}
  \begin{itemize}
    \item 새 데이터를 분류하려면 \textbf{학습도 모델도 없이}:
      가장 가까운 $k$개의 기존 데이터가 뭐라고 답했는지 보고
      다수결로 따른다
    \item 이 단순한 한 줄이 kNN이다 --- 거리를 \textbf{재는 데서
      출발}한다는 공통점이 k-means와 한 장에 묶는 이유
  \end{itemize}
  \vskip0.8em
  \begin{block}{kNN의 성격}
    \kb{게으른(lazy)}: 예측 시점에 거리를 재는 순간 모델이 완성된다 \\
    \kb{비모수(nonparametric)}: 파라미터 없이 점과 점을 비교
  \end{block}
  \kq{새 점의 정답은 ``그 점과 가까운 점들의 정답 평균''일까?}
\end{frame}
